\documentclass[11pt]{article}
\usepackage[margin=1.15in]{geometry}
\usepackage{amsmath,amssymb,amsthm}
\newtheorem{proposition}{Proposition}
\usepackage[T1]{fontenc}
\usepackage{setspace}
\usepackage{graphicx}
\usepackage[dvipsnames]{xcolor}
\usepackage[colorlinks=true,linkcolor=blue!50!black,citecolor=blue!50!black,urlcolor=blue!50!black]{hyperref}
\setstretch{1.08}
\setlength{\parskip}{0.55em}
\setlength{\parindent}{0pt}

\title{Syntactitude and Semantic Trajectories}
\author{Flyxion\\\small Independent Researcher}
\date{October 2026}

\begin{document}
\maketitle

\begin{abstract}
\noindent
The word \emph{syntactitude} appears once in the text of Daniel Dennett's \emph{Elbow Room}, on page 31 of the 1984 edition, and once in its index. Dennett introduces it as a synonym for a-rationality and reports that Douglas Hofstadter suggested it to him. This essay does three things. It fixes what the passage says and what it does not say, separating the term from the neighbouring notion of a syntactic engine and from the popular gloss that grammar is indifferent to morality. It examines a generated search overview that, given the exact term and the author's name, substituted the neighbouring notion and denied that Hofstadter coined the word without mentioning that Dennett reports Hofstadter suggested it. And it develops the term's relevance to a geometric account of meaning in which interpretation is a trajectory through a structured space, conceptual blending reshapes that space, and syntactitude is modelled as the possibility that a mechanical process comes to rest at an interpretation that reasons would not accept. The model is proposed and not tested empirically. Its formal results are modest: tangential capture is possible, a weak corrective intervention can move an unacceptable attractor without removing its inadequacy, and uniform guarantees fail when an observation does not distinguish contexts that require incompatible outcomes.
\end{abstract}

\section{A word on a page}

Some terms survive only in a single sentence of a single book. They are easy to remember in outline and easy to lose in detail: the topic, the author, the general atmosphere of the argument persist, while the exact word dissolves into a nearby one. \emph{Syntactitude} is such a term. A reader who remembers a discussion of grammar's indifference to meaning, and of Hofstadter's connection to it, is remembering something real. The word itself, however, is more specific than that memory suggests, and its specificity is the reason it is worth recovering.

The passage is on page 31 of \emph{Elbow Room} \cite{dennett1984}, in the chapter titled ``Making Reason Practical.'' The relevant sentence reads ``our a-rationality, or syntactitude'', followed by an aside that this is a term suggested to Dennett by Hofstadter. The index corroborates the occurrence and the page reference, though it is not testimony about provenance. In the 1984 pagination the entry reads ``syntactitude, 31'', adjacent to entries for \emph{syntactic engines} (28) and \emph{syntax} (28 to 29). A later reissue, consulted in electronic form, repaginates the book: there the sentence falls where the index entry says 35, and the index lists \emph{syntactic engines} at 32. Page references in this essay follow the 1984 edition. The term is therefore not an inference from a thematic resemblance. It is a named, indexed concept, and the only provenance documented here is Dennett's own report.

Dennett's wording introduces it as ``a-rationality, or syntactitude'', adding that the term was suggested to him by Hofstadter. Two cautions follow from the text as it stands. First, the attribution is Dennett's report, and what it supports is that Hofstadter suggested the word to Dennett. It does not by itself establish that no earlier use exists, and it says nothing about the circumstances in which the suggestion was made. Second, the word is introduced in passing, in a parenthetical gloss, inside an argument whose concerns lie elsewhere. Everything of interest in the term has to be recovered from that argument.

\section{What the passage says}

The argument of the surrounding pages runs as follows. Brains are, in one description, meaning manipulators, which Dennett calls semantic engines. In another description they are physical organs that can respond only to structural or formal properties of stimuli, which makes them syntactic engines. Syntax alone does not determine semantics, so a perfect semantic engine would be as impossible as a perpetual motion machine. Brains, Dennett concludes, only approximate the behaviour of the ideal, pure semantic engine. That engine, the perfect Kantian rational will, is ``friction-free, infinitely alert to nuances of meaning, perfectly invulnerable to sphexishness'' and physically impossible. Brains are ``very good substitutes'', cunningly designed to be ``not infinitely but indefinitely sensitive to meaningful changes, not perfectly but practically invulnerable to sphexishness.'' Sphexishness is the rut-bound behaviour illustrated by the digger wasp \emph{Sphex}, a term Dennett takes from Hofstadter.

The mechanism by which the approximation improves is reflexive. Dennett, following Hofstadter, says the explosive advance comes when pattern recognition is turned in upon itself. A creature sensitive to patterns in its own reactions can notice that it is caught in a futile loop or rut and leap out of it. One ``goes meta'' by representing one's representations, reflecting on one's reflections, reacting to one's reactions.

Page 31 then states the limit. A system that can reflect on its own activities in this way cannot reflect on all of them at once. It may reflect on its reflections, but sooner or later it exhausts its moves, its room, or its time. A finite, nonmagical perfect self-watcher is impossible, a point Dennett supports by appeal to the unsolvability of the halting problem as Hofstadter presents it. Anything short of perfection leaves a residue: in principle, the system could be placed in a situation that exposes its marginal and ordinarily invisible irrationality, its a-rationality, its syntactitude. Dennett insists that this should be neither a shock nor a disappointment, since nobody needs an argument from physical finitude to be persuaded that humans are not perfect meaning-extractors.

A note attached to the ``indefinitely sensitive'' sentence adds a detail that bears directly on the term. Dennett reports that Hofstadter makes excellent observations about our ``inevitable, finite susceptibility to absurd ruts'' of the kind Sphex falls into, and claims that our ruts, unlike hers, are recognizable in ``amiable quirks as artistic and cognitive styles'', moderately disabling on occasion but ``a small price to pay for our ability to think fast under the pressure of real time.'' Syntactitude, read alongside this note, is not merely a defect. It is the cost side of an economy in which speed and sensitivity are purchased by accepting a bounded vulnerability.

\section{What the term is not}

Two substitutions are tempting, and both change the meaning.

The first substitutes \emph{syntactic engine} for \emph{syntactitude}. A syntactic engine is any mechanism responding only to formal properties. Dennett applies the phrase to brains and to computers alike, as a description of what physical systems can do at the level of mechanism. Syntactitude belongs to a different descriptive level. It names a vulnerability in the performance of a mechanism that approximates a semantic engine: the leftover susceptibility that remains after the approximation has been made as good as finite resources allow. The two terms therefore do not divide mechanisms into exclusive kinds. Dennett's brains are syntactic engines, and they are also approximators of semantic engines whose performance shows syntactitude. A definition of syntactitude as the state of being a purely syntactic mechanism confuses the description of the mechanism with the description of the vulnerability in its performance.

The second substitution is the gloss that grammar is amoral, that a sentence can be well formed while being false, cruel, or absurd. The gloss is true and useful, and it lies close to a theme of the passage: form does not carry its own warrant. But it is a different point. It concerns the relation between well-formedness and truth or ethics. Dennett's concerns here are the relation between mechanism and meaning-sensitivity, and the limits of self-monitoring. A reader may well use the grammar gloss as a mnemonic, but it should be labelled as such rather than attributed to the text.

The syntactic-engine thesis also has critics, and the term has circulated since 1984. Labuschagne and Heidema \cite{labuschagne2005} contest the view of the brain as a syntactic engine limited to processing symbols by their structural properties, arguing for a dual-process model in which one mode of cognition is essentially semantical. Leon \cite{leon1990}, writing six years after \emph{Elbow Room}, uses the word once, in the opening paragraph on page 343. Dennett's challenge, he says, leads to two worries about reconciling mechanism and rationality. The second ``begins with mechanism and ends with the fear of our `syntactitude'{}; the fear that we are syntactic engines the functioning of which requires no reference to contentful states of the kind considered essential to (our) rationality.'' He then sets that worry aside: ``Here, it is the first worry that I will be examining.'' Leon cites Dennett 1984 for the challenge, does not mention Hofstadter, and does not gloss the word as a-rationality. His use postdates Dennett's and so does not bear on who suggested the word. It shows the term circulating as a name for a worry, which is a narrower thing than the residual susceptibility of an approximator that page 31 describes, and there is a case that this is itself a drift of the term. The full text of the paper was consulted, and the sentence quoted is the only occurrence.

Schollmeier \cite{schollmeier2009}, in a note on page 112, agrees ``with Dennett that we remain subject to what he dubs `syntactitude' (\emph{Room} 2.~31)'' but argues that ``our apparent syntactitude is due not to an imperfection in our rationality but to an imperfection in our practical application of it. There simply is not time enough or world enough.'' His phrasing establishes that Dennett uses the designation, and it need not assign him its invention; Dennett himself reports the term as suggested to him by Hofstadter. The two later uses show the term circulating after 1984, and neither is needed for the attribution, which rests on Dennett's own sentence.

A further trace sits in the reference list of a 2012 master's thesis at the University of the Witwatersrand, Andrew Nicholson's \emph{Six tricky pieces: naturalistic accounts of mental content and the problem of determinacy} \cite{nicholson2012}. A search excerpt from the repository copy shows an entry reading ``The Problem of Syntactitude. Unpublished manuscript.'', directly above ``Nicholson, A. (2012). Dretske and Fodor on Information. Unpublished manuscript.''. The excerpt is cut off before the author and year of the first entry, so the author is not established here; since the list is alphabetical and the next entry is Nicholson's, the entry may be his own, but that is an inference. The thesis text could not be retrieved, so it is not known where the manuscript is cited in the body, whether the thesis discusses it, or whether it uses the word in Dennett's sense. The search overview discussed below describes the term as associated with an unpublished paper on naturalistic accounts of mental content. Two later AI Mode answers to a query naming the thesis go further, and they disagree with each other. Both enumerate six ``tricky pieces'' and put the problem of syntactitude third, and both define it as the problem of individuating the vehicles of thought without appeal to meaning. Beyond that they diverge. The first lists the disjunction problem, distance, syntactitude, vacuous terms, intensionality, and abstract or mathematical content. The second lists the disjunction problem, the Swampman problem, syntactitude, distance, co-extensive properties, and a normativity gap, and adds an unsourced claim that the title is a nod to Feynman's \emph{Six Easy Pieces}. The first cites the thesis repository link in its source list; in the second the link appears only as a bare line at the end, attached to no claim. Two samples do not measure anything, but they show the pattern this essay is about: a stable core (the same position, the same definition) with the surrounding detail regenerated differently each time. A stable definition across samples shows that the generation is consistent, not that it is correct, and that definition is not Dennett's. None of the lists could be checked against the thesis text, so neither is used here as evidence of what the thesis contains. What the reference list does show is a manuscript whose title names a ``problem of syntactitude'', which is a further sign of the term circulating as a topic after 1984. The entry was seen only as an excerpt, and the manuscript itself was not seen.

Keeping these apart matters because each substitute makes the term sound more familiar than it is. The distinctive content is the combination of three claims: that competence at tracking meaning is real, that it is produced by mechanisms that do not track meaning as such, and that reflection improves the competence without ever completing it.

\section{A specimen: the search overview}
\label{sec:specimen}

A generated overview offers a convenient specimen of how that distinctive content is lost. A search was made on the exact term in quotation marks together with the author's surname. The overview returned in two successive screenshots. The first begins by saying that Hofstadter did not coin the word, then says that it directly addresses the philosophical anxiety at the centre of his work on artificial intelligence and consciousness, and cites an encyclopedia of science fiction as its source. The second defines the term as a philosophical word for being a purely syntactic mechanism, a system that processes symbols and rules without inherent understanding of their meaning or content; in the first capture the same definition begins ``the condition of being a purely syntactic engine''. Both captures are reproduced in Figure~\ref{fig:overview}.

\begin{figure}[ht]
\centering
\includegraphics[width=0.30\textwidth]{overview1.jpg}\hspace{2em}\includegraphics[width=0.30\textwidth]{overview2.jpg}
\caption{Captures of the search overview for the query ``syntactitude'' Hofstadter, taken on the same day (left: 6:26; right: 6:27). The left capture contains the statement that Hofstadter did not coin the word and the SF Encyclopedia citation. Both contain the definition in terms of a purely syntactic engine or mechanism.}
\label{fig:overview}
\end{figure}

Set against the page, the overview fails in three ways, and the failures are of different kinds.

\begin{enumerate}
\item \textbf{Substitution.} The definition given is the definition of a syntactic engine. It is the neighbouring concept from page 28, delivered in place of the one on page 31.
\item \textbf{Attribution.} The overview says Hofstadter did not coin the word. A denial of coinage would not strictly contradict Dennett's report that Hofstadter suggested the term to him, and the page does not establish coinage either. The failure is one of omission and sourcing: the overview offers a flat denial without mentioning the one primary record that bears on the question, and the source it cites is not that book.
\item \textbf{Generalization.} The claim that the word addresses the anxiety at the heart of Hofstadter's work converts a specific usage into a thematic gesture. It is the kind of sentence that is difficult to falsify because it is not about anything in particular.
\end{enumerate}

A further capture, in Google's separate AI Mode and supplied here as pasted text without a screenshot, answers the query ``syntactitude'' alone. It defines the term as ``the condition of being a purely syntactic engine'' and as ``the fear or philosophical premise'' that minds operate mechanically without any inherent connection to meaning, and its listed sources include a copy of Leon's ``The Mechanics of Rationality'' \cite{leon1990}. That definition tracks Leon's phrasing about the fear of being syntactic engines more closely than it tracks page 31 of \emph{Elbow Room}. If that is where it comes from, the substitution has a traceable route: a later author's use of the word as a name for a worry, which may itself have drifted from Dennett's sense, was taken as the definition. Leon's text has been checked: the word occurs there once, as the name of a fear that he then sets aside. What AI Mode drew on beyond that is not known, so the route remains an inference from the match in wording.

The result is fluent and confident, and it recognizes the word. That recognition is what makes it dangerous: a reader who sees the term acknowledged, defined, and placed in a plausible intellectual context has little reason to suspect that the definition belongs to another term. Only a return to the book, and to its index, restores the distinction. The lesson is not that summaries are unreliable in general. It is that a summary can be unreliable in exactly the dimension that fluency conceals, and that primary text with an index remains a better instrument for questions of the form ``where, precisely, does this word occur and what does it say.''

The episode is also, quite directly, an instance of the thing it concerns. A system that responds to the surface form of a query, the quoted term and the name, and produces a continuation shaped by nearby associations, will produce a coherent answer whether or not the associations preserve the relevant distinction. Whether this is an instance of the capture modelled later, or some other failure such as a defect of the embedding, is a question Section~\ref{sec:relation} takes up, and it is not settled by the resemblance alone.

\section{Contextual descent and interpretive adequacy}

The connection to a geometric account of meaning arises from asking how mechanisms with no intrinsic concern for significance can nonetheless produce behaviour that is sensitive to it. Conceptual blending, as developed by Fauconnier and Turner \cite{fauconnier2002}, supplies one part of an answer: structure from two input domains is selectively projected into a blended space, where it is composed, completed, and elaborated, yielding inferences that neither input supplied alone. The coupled-potential model below is a model of blending offered here, and is not a transcription of Fauconnier and Turner's theory, which contains no such formalism. The sections that follow state a minimal model in which that productivity, and the failure that syntactitude names, can both be made precise.

Let $(M,g)$ be a smooth, compact Riemannian manifold without boundary, representing the admissible interpretive states of a specified task. Compactness is a simplifying assumption that ensures completeness of the smooth dynamics considered below. For each context $c\in C$, let
\[
V_c\colon M\longrightarrow\mathbb{R}
\]
be a twice continuously differentiable potential. The intrinsic descent dynamics are
\[
\dot{x}(t)=-\operatorname{grad}_g V_c(x(t)).
\]
Perturbations, where they are considered, are taken to be tangent vectors, $\eta(t)\in T_{x(t)}M$. The metric is part of the model: it determines how differences between interpretations are measured and how a differential of the potential becomes a direction of movement. A potential alone does not determine a trajectory.

Along any solution,
\[
\frac{d}{dt}V_c(x(t))
=
dV_c\bigl(\dot{x}(t)\bigr)
=
-\left\|\operatorname{grad}_g V_c(x(t))\right\|_g^2
\leq 0.
\]
Descent therefore supplies a certificate of decreasing potential. It supplies no certificate that the potential expresses the relevant reasons adequately. To represent that further condition, let $R_c\subseteq M$ be a nonempty closed set of acceptable interpretations, specified independently of the descent dynamics.

Suppose $m$ is a critical point of $V_c$ whose Riemannian Hessian is positive definite. Write $H_c$ for the self-adjoint operator on $T_mM$ defined by $g(H_cv,w)=\operatorname{Hess}_g V_c(m)[v,w]$, which is identified with a symmetric matrix in normal coordinates, so that the condition reads:
\[
\operatorname{grad}_g V_c(m)=0,
\qquad
\operatorname{Hess}_g V_c(m)[v,v]>0
\quad
\text{for every }0\neq v\in T_mM.
\]
Then $m$ is a strict local minimum and a locally asymptotically stable equilibrium of gradient descent. If $m\notin R_c$, its stability constitutes interpretive capture: nearby trajectories converge toward an unacceptable interpretation while remaining entirely within the admissible state space.

Confinement and adequacy are thus logically independent. An explicit example takes $M=S^1$ with its standard metric, and
\[
V(\theta)=1-\cos\theta,
\qquad
R=\{\theta:d_{S^1}(\theta,\pi)\leq\delta\},
\qquad
0<\delta<\frac{\pi}{2}.
\]
The dynamics satisfy $\dot{\theta}=-\sin\theta$. Every initial state except $\theta=\pi$ converges to $\theta=0$, which lies outside $R$. No trajectory leaves the circle. The failure is produced by the relationship between the potential and the acceptability set, and not by a departure from the manifold.

This construction establishes that tangential capture is possible. It does not establish that every finite interpretive system must exhibit it on every task. A particular task can have a potential whose minima all lie in $R_c$, and nothing in Dennett's argument denies this. Dennett invokes the halting problem to motivate a limit on unrestricted self-inspection. This does not by itself establish a general theorem of interpretive inadequacy. The obstruction in Section~\ref{sec:ambiguity} rests instead on an explicit informational condition: indistinguishable contexts have incompatible acceptable outcomes.

\section{Basins, convergence rates, and capture risk}

Write $\phi_t^c$ for the flow induced by $V_c$. For an attracting minimum $m$, define its basin by
\[
B_c(m)
=
\left\{
x\in M:
\lim_{t\to\infty}\phi_t^c(x)=m
\right\}.
\]
Let $\mu_c$ be a probability measure describing the distribution of initial interpretive states. If the minima are isolated, a task-relative measure of capture is
\[
\mathcal{S}(c;\mu_c)
=
\mu_c\left(
\bigcup_{\substack{m\in\operatorname{Min}(V_c)\\m\notin R_c}}
B_c(m)
\right).
\]
The quantity is introduced as a definition and is not applied in what follows. It measures the probability of convergence to an unacceptable minimum under the specified initialization. It does not count nonconvergent trajectories or other failure modes, and it is not an intrinsic scalar property of the mechanism. Changing the context, the adequacy criterion, or the distribution of initial states can change it.

The speed of local convergence depends on curvature. In normal coordinates $z$ around a nondegenerate minimum $m$,
\[
V_c(z)
=
V_c(m)+\frac12 z^\top H_c z+o(\|z\|^2),
\]
where $H_c$ is positive definite, the remainders being those available for $V_c\in C^2$. Riemannian normal coordinates at $m$ give $g_{ij}(m)=\delta_{ij}$ and $\partial_kg_{ij}(m)=0$, so the connection terms vanish at $m$ and the linearization is the Euclidean one. Consequently,
\[
\dot{z}=-H_cz+o(\|z\|).
\]
The linearized dynamics give
\[
\|z(t)\|
\leq
e^{-\lambda_{\min}(H_c)t}\|z(0)\|.
\]
A larger smallest Hessian eigenvalue produces faster local contraction in this approximation. Basin depth and basin volume do not alone determine that rate, and global convergence behaviour also depends on gradients, mobility, perturbations, and initial conditions.

The distinction matters for the economy of interpretation suggested by Dennett's note on Hofstadter. A familiar basin can attract a large proportion of initial states, contract rapidly near its minimum, and resist perturbation. These are separate properties. Strong attraction can facilitate rapid settling, which is valuable under real-time pressure, while making redirection difficult. Fluency may accompany rapid convergence without supplying evidence that the attracting minimum belongs to $R_c$. Settling alone supplies no certificate of correctness, although a system may of course detect its own error afterwards.

\section{Meaning as a path-dependent functional}

An endpoint account assigns significance through a function $f\colon M\to\mathbb{R}$. A trajectory account permits a functional on paths,
\[
\mathcal{I}_c\colon\mathcal{P}(M)\longrightarrow\mathbb{R},
\]
where $\mathcal{P}(M)$ is a chosen class of piecewise smooth interpretive trajectories. One elementary construction uses a smooth one-form $\omega_c$:
\[
\mathcal{I}_c[\gamma]
=
\int_\gamma\omega_c
=
\int_0^T
\omega_c\bigl(\gamma(t)\bigr)
\bigl[\dot{\gamma}(t)\bigr]\,dt.
\]
This functional is invariant under orientation-preserving reparameterization. It is therefore insensitive to the speed at which the route is traversed, including pauses, and records only the route.

If $\omega_c=df_c$ is exact, then
\[
\mathcal{I}_c[\gamma]
=
f_c(\gamma(T))-f_c(\gamma(0)),
\]
so the functional reduces to endpoint information. Genuine path dependence requires failure of exactness. For two paths $\gamma_1,\gamma_2$ with the same initial and terminal states,
\[
\mathcal{I}_c[\gamma_1]
-
\mathcal{I}_c[\gamma_2]
=
\oint_{\gamma_1*\overline{\gamma_2}}\omega_c.
\]
When the resulting loop bounds an oriented surface $\Sigma$, Stokes' theorem gives
\[
\mathcal{I}_c[\gamma_1]
-
\mathcal{I}_c[\gamma_2]
=
\int_\Sigma d\omega_c.
\]
Thus $d\omega_c$ measures local route dependence in this model. A closed but nonexact one-form can additionally retain dependence on the global topology of the interpretive space.

For example, on $\mathbb{R}^2$ take $\omega=u\,dv$. A path from $(0,0)$ to $(1,1)$ that first changes $u$ and then changes $v$ has integral $1$. Reversing that order gives integral $0$. The endpoint is identical, but the ordered development differs. These two paths are not asserted to be solutions of the same autonomous gradient system with the same initial condition. They represent different contextual histories or interventions. For a fixed context and initial state, the $C^1$ gradient vector field determines a unique trajectory. The construction exhibits one mathematically explicit form of path dependence. It does not assert that every path-dependent interpretation admits representation by a one-form integral: dependence on elapsed time, accumulated memory, or other features is excluded by this representation.

One reading takes $d\omega_c$ as the local density of context that changes the significance of a step, so that a step has different weight depending on what has already been traversed. That reading is offered as an interpretation of the construction and is not derived from it. The integral is not proposed as an exhaustive measure of meaning. It exhibits a precise way in which an interpretive history can carry information that its endpoint does not preserve. Sequences of association, evidential acquisition, and correction can be represented by additional functionals of the same trajectory.

\section{Conceptual blending as coupled constraint formation}

Let $(M_1,g_1)$ and $(M_2,g_2)$ represent two input domains. Suppose each domain carries a potential $V_i$, and let
\[
q_i\colon M_i\longrightarrow Z
\]
map domain-specific states into a shared Euclidean feature space $Z$. The maps encode selected correspondences rather than a complete identification of the domains.

On $M_1\times M_2$, equipped with the product metric, define
\[
W_\lambda(x,y)
=
V_1(x)+V_2(y)
+
\frac{\lambda}{2}
\|q_1(x)-q_2(y)\|^2,
\qquad
\lambda\geq 0.
\]
The final term rewards compatibility of the projected features. Writing $r=q_1(x)-q_2(y)$, its critical-point equations are
\[
\operatorname{grad}_{g_1}V_1(x)
+
\lambda(Dq_1)_x^*r
=0,
\qquad
\operatorname{grad}_{g_2}V_2(y)
-
\lambda(Dq_2)_y^*r
=0,
\]
where the adjoints are defined by the relevant metrics.

At $\lambda=0$, the domains evolve independently. For positive coupling, equilibria depend jointly on both domains and need not coincide with either domain's uncoupled minima. This is a minimal mathematical representation of how bringing domains into relation can change the available stable interpretations.

The coupled system can also be reduced. Assume $V_i,q_i\in C^2$, so that $W_\lambda\in C^2$. If, for each $x$, there is a unique minimizing state $y_*(x)$ that varies as a $C^1$ branch, define
\[
V_{\mathrm{eff}}(x)=W_\lambda(x,y_*(x)).
\]
In local coordinates, assume
\[
W_{yy}\succ 0,
\]
which gives, through the implicit function theorem, a local $C^1$ branch of strict local minima through any existing stationary point of $W_\lambda(x,\cdot)$. It does not by itself establish that this branch is the global minimizer or that it is unique; those properties are supplied by the preceding assumption on $y_*(x)$. Then
\[
D^2V_{\mathrm{eff}}
=
W_{xx}-W_{xy}W_{yy}^{-1}W_{yx},
\]
evaluated at $(x,y_*(x))$. Here $D^2$ denotes ordinary second derivatives in local coordinates. At a critical point of $V_{\mathrm{eff}}$ these represent the intrinsic Hessian, while away from critical points the Riemannian Hessian also includes connection terms. The cross-domain terms change the curvature of the effective landscape. Coupling can therefore alter stability even when the separate input potentials remain fixed. A supplied example is one possible means of changing such coupling constraints. The equations do not establish a general division between examples and instructions, since either can alter the effective context and landscape.

This construction models selected correspondence and joint elaboration. It does not exhaust conceptual blending, whose selective projection, completion, and emergent organization may require richer state spaces and additional constraints.

\section{Reflection and the persistence of equilibria}

Represent a reflective intervention by a perturbation
\[
V_{c,\varepsilon}=V_c+\varepsilon P_c,
\]
where $P_c\in C^2(M)$ penalizes a recognized pattern. Suppose $m$ is a nondegenerate minimum of $V_c$. Applying the implicit function theorem to
\[
\operatorname{grad}_g(V_c+\varepsilon P_c)(m_\varepsilon)=0
\]
produces a nearby critical point $m_\varepsilon$ for sufficiently small $\varepsilon$. Its first-order displacement is
\[
\left.\frac{d m_\varepsilon}{d\varepsilon}\right|_{\varepsilon=0}
=
-H_c^{-1}\operatorname{grad}_g P_c(m),
\]
where $H_c$ is the Hessian operator of $V_c$ at $m$. Positive definiteness persists under sufficiently small perturbations, so the nearby equilibrium remains attracting.

If $m\notin R_c$ and $R_c$ is closed, then $d_g(m,R_c)>0$. Continuity implies that sufficiently small perturbations leave $m_\varepsilon$ outside $R_c$. A weak reflective intervention can consequently move an unacceptable attractor without eliminating its inadequacy.

This is a local persistence result, not a theorem that correction must fail. A stronger intervention can move an equilibrium into $R_c$, destroy it through a bifurcation, or change the state space and context themselves. Reflection may also alter the metric or the dynamics rather than merely adding a potential term.

\begin{proposition}[Persistence of capture under weak correction]
Let $V_c$ and $P_c$ be $C^2$, let $m$ be a nondegenerate minimum of $V_c$ with $m\notin R_c$, and let $R_c$ be closed. For all sufficiently small $\varepsilon$, the perturbed potential $V_c+\varepsilon P_c$ has a nondegenerate minimum $m_\varepsilon$ near $m$, and $m_\varepsilon\notin R_c$.
\end{proposition}

The mathematical connection to syntactitude is therefore precise but limited: noticing a failure and introducing a corrective influence does not by itself guarantee its removal. Successful correction depends on the intervention's effect on the relevant attractor and on the criterion of adequacy.

\section{Contextual ambiguity and limits on guarantees}
\label{sec:ambiguity}

A separate obstruction arises when the available observation does not distinguish contexts that require incompatible outcomes. Let
\[
h\colon C\longrightarrow O
\]
describe the observation supplied to an interpretive system. A deterministic endpoint policy is a map $a\colon O\to M$. It succeeds uniformly when
\[
a(h(c))\in R_c
\qquad
\text{for every }c\in C.
\]

For an observation $o$, define the common acceptable set
\[
K(o)
=
\bigcap_{c\in h^{-1}(o)}R_c.
\]
A uniformly successful deterministic policy exists precisely when one can select an element of $K(o)$ for every observation in $h(C)$. In particular, if $K(o)=\varnothing$, uniform success is impossible at that observation. The proof is immediate: the system produces the same endpoint for every context sharing $o$, and that endpoint would have to belong to every corresponding acceptability set.

This obstruction is independent of convergence speed and computational power. A reflective procedure receiving no additional contextual information remains subject to it. An additional observation $q(c)$ refines the context class to
\[
\{c:h(c)=o,\ q(c)=z\},
\]
with common acceptable set
\[
K(o,z)=\bigcap_{c:\,h(c)=o,\ q(c)=z}R_c\ \supseteq\ K(o),
\]
the inclusion holding because the index set of $K(o,z)$ is a subset of the index set of $K(o)$, and an intersection over a subset can only be larger.

Returning to a primary source can function as such a refinement. It supplies information that a term and an author's name may leave unresolved. The improvement arises from changing what the system can distinguish, rather than merely increasing the force with which it pursues an interpretation.

There is also a finite probabilistic version. Suppose contexts $c_1,\ldots,c_k$ share an observation and their acceptable sets are pairwise disjoint. For any probability distribution $\nu_o$ over endpoints,
\[
\sum_{i=1}^k\nu_o(R_{c_i})\leq 1,
\qquad
\min_i\nu_o(R_{c_i})\leq\frac1k.
\]
The bound is tight: when the sets are disjoint, a uniform distribution over one chosen point from each attains $1/k$ for every context. Randomization therefore cannot provide uniformly high success across these incompatible contexts.

The obstruction is not a formalization of Dennett's argument. His concerns the limits of a system's inspection of its own operations, whereas this one concerns a coarse observation, and it can arise for a system that inspects everything available to it. It identifies a condition his discussion does not consider, available in the same conceptual space. These results do not prove universal failure of interpretation. They identify explicit conditions under which a guarantee is unavailable. Conversely, a sufficiently informative observation and an appropriately aligned potential can support reliable success on a specified task.

\section{Relation to manifold-aligned dynamics}
\label{sec:relation}

The space $M$ used above corresponds to the semantic manifold of the author's earlier work on manifold-aligned generative inference \cite{flyxion2025magi,flyxion2026hallucination}. In that work meaningful states occupy a constrained submanifold $M\subset\mathbb{R}^n$ of intrinsic dimension $d\ll n$, with the ambient space splitting at each point as $\mathbb{R}^n = T_xM\oplus N_xM$. Explanation is tangent and hallucination is normal: an update with nonzero normal component leaves $M$ and produces structure that no lawful regularity supports. In the later formulation $M$ is not posited but emerges as the maximal connected admissible subset of the space of histories. When $M$ carries the metric induced from the ambient space, it is an instance of the Riemannian manifold $(M,g)$ used here.

Reading syntactitude into that framework sharpens a distinction the earlier papers leave implicit. Two different failures are available.

\begin{enumerate}
\item \textbf{Normal drift.} The update acquires a component in $N_xM$ and the trajectory leaves the manifold. The remedy is projection onto $T_xM$, or an intrinsic gradient. Projection guarantees confinement only under suitable continuous dynamics. A finite tangent step can leave a curved manifold, so a discrete implementation needs a retraction or another constraint-preserving update.
\item \textbf{Tangential capture.} The dynamics remain on $M$ throughout, so no normal component appears, and the trajectory nonetheless converges to a minimum of $V_c$ outside $R_c$. Section 5 gives an explicit example. The failure lies in the relation between the potential and the criterion of adequacy.
\end{enumerate}

The remedy for the first failure does not touch the second, because confinement and adequacy are independent. Tangency is therefore a condition on where a process may go and is silent on whether it ends where it should. This is a consequence of the model as stated here, offered as a refinement of the earlier framework. It is not a claim the earlier papers make, and no experiment in this essay tests it.

A related observation from the later paper is that many semantic pathologies are failures of the embedding that places an intrinsically coherent trajectory in ambient coordinates, and are not failures of $M$ itself. Capture may combine with this: a basin may be a genuine feature of $M$ that is nevertheless wrong for a given context, because the context was under-specified in the prompt, the query, or the interface.

The sheaf-theoretic layer of the earlier framework can be applied to the specimen of Section~\ref{sec:specimen} in a toy case small enough to check by hand. Let $A=\{\delta,\ell\}$ be the attributes about which sources speak: $\delta$, the definition of the term, and $\ell$, its location in the book. For each subset $U\subseteq A$ let $S(U)$ be the set of assignments of a value to each attribute in $U$, with restriction maps given by restricting assignments. This is the sheaf of functions on the finite set $A$, so compatible local sections over a cover glue uniquely. Three sources supply local sections:
\begin{itemize}
\item the page supplies a section $s_P$ on the whole of $A$, with $s_P(\delta)=d_P$, the definition of a-rationality as the residual susceptibility of an approximator to a semantic engine, and $s_P(\ell)=31$ in the 1984 pagination;
\item the index: $s_I$ on $\{\ell\}$, with $s_I(\ell)=31$;
\item the overview: $s_O$ on $\{\delta\}$, with $s_O(\delta)=d_O$, the definition as being a purely syntactic mechanism.
\end{itemize}
On the overlap $\{\ell\}$ of the page and the index the sections agree. On the overlap of the index and the overview, which is empty, they agree trivially. On the overlap $\{\delta\}$ of the page and the overview they disagree, because $d_P\neq d_O$. The family $\{s_P,s_I,s_O\}$ is therefore not compatible, and no section over $A$ restricts to both $s_P$ and $s_O$. The compatible subfamily $\{s_P,s_I\}$ glues uniquely to $s_P$. The sheaf axiom is not violated, since it governs only compatible families. What the example exhibits is the location of the incompatibility: it sits on a named overlap, the definition, and nowhere else.

The example has limits. It models substitution only: the failure modes of omitting a source's report and of generalizing a specific usage into a theme are not encoded in it. The judgment that $d_P\neq d_O$ is the one argued for in Sections~3 and~4, and the toy encodes it without deriving it. Attribution is deliberately absent from $A$, because a denial of coinage and a report of suggestion are not assigned incompatible values by anything shown here. And the example concerns agreement among sources. It does not show that the overview's claims belong to $M$, and so it does not show that the overview exhibits tangential capture and not some other failure. That question has not been established here. The evidence available shows a substitution of a neighbouring concept and a denial of coinage that omits the primary source's report of Hofstadter's suggestion, which is a narrower finding.

\section{Working with language models}

A language model offers an unusually clean setting in which to observe these dynamics, because the conditions that shape its continuations are largely explicit. A prompt supplies a context under which some continuations become easier than others. Examples alter the landscape in the manner of blending. Within the proposed model, a correction may be represented by a penalty on the preceding pattern or by a change of context; neither representation is established here as the internal mechanism of an actual language model. The fluent but inadequate answer is the characteristic failure: a continuation that has settled into a basin supported by familiarity and surface cues, without the distinction that the task required.

This suggests a practical division of labour. The model supplies fluency, breadth of association, and the capacity to elaborate a blend rapidly. The human contribution is partly a trained attention to ambiguity, presupposition, and the difference between what an assertion says and what it takes for granted. Those are the skills that detect when a trajectory has settled in the wrong place, and that identify which example, distinction, or constraint would reshape the landscape so that a better basin becomes reachable. A background in linguistics is one source of such attention, though not the only one. It does not confer knowledge of a model's internals.

The specimen of section four illustrates the remaining practice. When the question is where a word occurs and what it says, the most reliable intervention is to leave the fluent channel and consult the source, and where possible the index, which is a structure built for exactly that question and indifferent to plausibility.

\section{Conclusion}

Dennett's single sentence does a surprising amount of work. It names a vulnerability that competence cannot abolish, distinguishes it from the mechanism-level fact that brains are syntactic engines, ties it to a formal limit on self-inspection, and, through a note, treats it as the cost side of a bargain that buys speed. Recovering the word accurately, rather than a nearby one, preserves each of those distinctions. The generated overview that replaced it with the neighbouring definition shows how easily they are lost.

The geometric reading developed here is an extension rather than an exegesis. It treats interpretation as a trajectory, blending as a change in the constraints under which trajectories run, and syntactitude as the possibility that descent settles at an interpretation outside the acceptable set while remaining entirely within the admissible space. The formal results are modest and stated with their conditions: capture is possible, weak correction can leave it in place, and guarantees fail where an observation does not distinguish contexts that require incompatible outcomes. Whether the model earns its keep is an empirical matter. What it already provides is a vocabulary in which reformulating, exemplifying, and returning to the source can be described together, as ways of changing the landscape or the available information. Such changes support iterative improvement and checking against sources. They do not guarantee that the process arrives where reasons would have it rest.

\begin{thebibliography}{9}
\bibitem{flyxion2025magi}
Flyxion, \emph{Never Predict Noise: Manifold-Aligned Prediction as Generative Inference}. November 2025, revised May 2026. \url{https://standardgalactic.github.io/alphabet/Never-Predict-Noise.pdf}.

\bibitem{flyxion2026hallucination}
Flyxion, \emph{Hallucination Is Normal: A Geometric Theory of Manifold-Aligned Semantic Dynamics}. February 18, 2026. \url{https://standardgalactic.github.io/library/hallucination-is-normal.pdf}.

\bibitem{labuschagne2005} W. A. Labuschagne and J. Heidema, ``Natural and artificial cognition: On the proper place of reason.'' \emph{South African Journal of Philosophy} 24(2): 137--151, 2005. doi:10.4314/sajpem.v24i2.31421.
\bibitem{leon1990} Mark Leon, ``The Mechanics of Rationality.'' \emph{The Southern Journal of Philosophy} 28(3): 343--366, 1990 (pages as numbered in the copy consulted; the last page holds the reference list). The word occurs once, on p.\ 343.
\bibitem{schollmeier2009} Paul Schollmeier, \emph{Human Goodness: Pragmatic Variations on Platonic Themes}. Cambridge University Press, 2009. Chapter 3, ``Human Happiness'', p.\ 112, note 54. doi:10.1017/CBO9780511498688.005.
\bibitem{nicholson2012} Andrew Nicholson, \emph{Six tricky pieces: naturalistic accounts of mental content and the problem of determinacy}. M.A. thesis, University of the Witwatersrand, 2012 (published 9 October 2012). Seen only as a search excerpt of the reference list on 4 October 2026; full text not retrieved.
\bibitem{dennett1984} Daniel C. Dennett, \emph{Elbow Room: The Varieties of Free Will Worth Wanting}. MIT Press, 1984. Chapter 2, ``Making Reason Practical''; passage on p.\ 31 (1984 pagination; the index lists p.\ 35 in a later reissue); index entries for \emph{syntactic engines}, \emph{syntactitude}, and \emph{syntax}.
\bibitem{hofstadter1982} Douglas R. Hofstadter, ``Can Creativity Be Mechanized?'' 1982. Cited by Dennett as 1982b, pp.\ 27 to 28, for the unsolvability of the halting problem; the same paper is cited on p.\ 11 for the term sphexishness.
\bibitem{fauconnier2002} Gilles Fauconnier and Mark Turner, \emph{The Way We Think: Conceptual Blending and the Mind's Hidden Complexities}. Basic Books, 2002.
\end{thebibliography}

\end{document}

